\subsection{End-to-End Localization Example}\label{sec:end-to-end-dataloc}
\begingroup
\tcbset{datalocstep/.style={
    colback=gray!2, colframe=cyan!40!black, boxrule=0.5pt,
    boxed title style={colback=cyan!6, colframe=cyan!40!black},
    left=2mm, right=2mm, top=3mm, bottom=1.5mm,
    before skip=8pt, after skip=8pt}}
\setminted{fontsize=\footnotesize,baselinestretch=1,breaklines=true,breakautoindent=true}
\setminted[prolog]{xleftmargin=1.5em,numbersep=5pt}
Using Astropy commit \mbox{\texttt{d16bfe05a744909de4b27f5875fe0d4ed41ce607}}, this example traces a recorded localization interaction through the corresponding source code and extracted program facts.

\paragraph*{Original user query.}
The user submits the following localization request:
\begin{quote}
Find all classes where: (1) the \texttt{\_\_init\_\_} method contains an if statement that checks whether a parameter is None, (2) one branch of this if statement raises TypeError and the other branch raises ValueError
\end{quote}

\paragraph*{Relevant facts extracted from the codebase.}
\fref{fig:dataloc-extracted-facts}\subref{fig:dataloc-facts-timeseries} and \subref{fig:dataloc-facts-wcs} show relevant program facts for \texttt{TimeSeries} and \texttt{WCS}, respectively, following the schema in Table~\ref{tab:dataloc-fact-schema} and its accompanying description. These facts record the line ranges of \texttt{\_\_init\_\_} methods and conditionals, together with source lines containing explicit \texttt{raise} statements. 

\begin{landscape}
\begin{figure}[p]
\centering
\begin{subfigure}{\linewidth}
\begin{stepbox}[datalocstep]{Extracted facts: TimeSeries}
\begin{minted}[linenos=false,xleftmargin=0pt]{prolog}
function_definition("astropy/timeseries/sampled.py", "TimeSeries.__init__", 61, 135, 7, False, "TimeSeries").
control_structure("astropy/timeseries/sampled.py", "if_statement", 107, 130, "time_start IsNot None").
source_line("astropy/timeseries/sampled.py", 113, "raise TypeError(\\"'time' is scalar, so 'time_delta' is required\\")").
source_line("astropy/timeseries/sampled.py", 125, "raise ValueError(\\"Length of 'time' ({}) should match \\"").
\end{minted}
\end{stepbox}
\caption{Extracted facts for \texttt{TimeSeries}.}
\label{fig:dataloc-facts-timeseries}
\end{subfigure}
\par\medskip
\begin{subfigure}{\linewidth}
\begin{stepbox}[datalocstep]{Extracted facts: WCS}
\begin{minted}[linenos=false,xleftmargin=0pt]{prolog}
function_definition("astropy/wcs/wcs.py", "WCS.__init__", 380, 546, 12, False, "WCS").
control_structure("astropy/wcs/wcs.py", "if_statement", 391, 527, "header Is None").
source_line("astropy/wcs/wcs.py", 429, "raise TypeError(").
source_line("astropy/wcs/wcs.py", 413, "raise ValueError(").
\end{minted}
\end{stepbox}
\caption{Extracted facts for \texttt{WCS}.}
\label{fig:dataloc-facts-wcs}
\end{subfigure}
\par\medskip
\caption{Relevant program facts from the Astropy repository.}
\label{fig:dataloc-extracted-facts}
\end{figure}
\end{landscape}

\paragraph*{Initially generated Datalog query.}
\fref{fig:dataloc-query-repair}\subref{fig:dataloc-query-initial} shows the initial query generated by the agent. \texttt{InitFunc} targets \texttt{\_\_init\_\_} methods, and \texttt{IfNone} targets enclosed conditionals that test for \texttt{None}. The variable pairs \texttt{f\_start}/\texttt{f\_end} and \texttt{if\_start}/\texttt{if\_end} denote their inclusive line bounds. Argument types are specified in \sref{sec:fact-extraction-dataloc}; \texttt{\_} denotes an anonymous variable.

\paragraph*{Checking, repair, and refinement.}
Exact matching against \texttt{\_\_init\_\_} yields no \texttt{InitFunc} tuples because method names are class-qualified. The empty result triggers a Contains-Literal diagnostic mutation (\sref{sec:method:int-dataloc}), which replaces exact matching of \texttt{\_\_init\_\_} with substring matching. The mutated query produces a nonempty result, yielding \emph{fragile-empty} feedback that identifies the name constraint as a possible cause. Guided by this feedback, the agent inspects the recorded method names and revises the query to match the substring \texttt{.\_\_init\_\_}.

The next candidate query introduces \texttt{RaiseT} and \texttt{RaiseV} for lines containing \texttt{TypeError} and \texttt{ValueError} raises, and \texttt{HasBoth} for conditional spans containing both. It retrieves \texttt{TimeSeries} and \texttt{WCS}. However, finding both exception types within the same conditional span does not establish that they occur in different branches.

\fref{fig:dataloc-query-repair}\subref{fig:dataloc-query-repaired} presents the final query after further agent revisions. Here, \texttt{ST} denotes the recorded control-structure type, \texttt{c} the source-line context, and \texttt{l1} and \texttt{l2} the line numbers of the matched \texttt{TypeError} and \texttt{ValueError} raises, respectively. Compared with the preceding candidate query, it accepts either \texttt{if\_statement} or \texttt{unknown}, requires only the conditional's starting line to fall within the \texttt{\_\_init\_\_} method, and changes the lower-bound checks on raise locations from \texttt{>} to \texttt{>=}. These changes relax the candidate-selection constraints but return the same two candidates. Source verification is therefore still required to determine whether the raises occur in different branches.

Both panels retain the submitted \texttt{contains} argument order, corrected to Souffl\'e's substring-first convention by semantic validation before execution (\sref{sec:method:val-dataloc}). The figure omits only the shared fact-input include directive.

\begin{landscape}
\begin{figure}[p]
\tcbset{datalocstep/.append style={before skip=0pt,after skip=0pt,top=1mm,bottom=1mm}}
\setstretch{1}
\captionsetup{font=small,skip=3pt}
\centering
\begin{subfigure}[t]{\linewidth}
\vspace{0pt}
\centering
\begin{stepbox}[datalocstep]{Initially generated query}
\begin{minted}[linenos=true,fontsize=\fontsize{9}{10}\selectfont,baselinestretch=0.9]{prolog}
.decl InitFunc(file_path: symbol, class_name: symbol, f_start: number, f_end: number)
InitFunc(file_path, containing_class, f_start, f_end) :-
    function_definition(file_path, "__init__", f_start, f_end, _, _, containing_class),
    containing_class != "".
.decl IfNone(file_path: symbol, class_name: symbol, if_start: number, if_end: number, condition: symbol)
IfNone(file_path, class_name, if_start, if_end, condition) :-
    InitFunc(file_path, class_name, f_start, f_end),
    control_structure(file_path, "if_statement", if_start, if_end, condition),
    contains(condition, "None"),
    if_start >= f_start, if_end <= f_end.
.output IfNone
.output InitFunc
\end{minted}
\end{stepbox}
\caption{Initial query for \texttt{\_\_init\_\_} methods and enclosed conditionals.}
\label{fig:dataloc-query-initial}
\end{subfigure}
\par
\begin{subfigure}[t]{\linewidth}
\vspace{0pt}
\centering
\begin{stepbox}[datalocstep]{Final query submitted by the agent}
\begin{minted}[linenos=true,fontsize=\fontsize{9}{10}\selectfont,baselinestretch=0.9,highlightlines={3,4,8,9,10,11,12,13,14,15,16,17,18,19,20},highlightcolor=cyan!8]{prolog}
.decl InitFunc(file_path: symbol, class_name: symbol, f_start: number, f_end: number)
InitFunc(file_path, containing_class, f_start, f_end) :-
    function_definition(file_path, function_name, f_start, f_end, _, _, containing_class),
    containing_class != "", contains(function_name, ".__init__").
.decl IfNone(file_path: symbol, class_name: symbol, if_start: number, if_end: number, condition: symbol)
IfNone(file_path, class_name, if_start, if_end, condition) :-
    InitFunc(file_path, class_name, f_start, f_end),
    control_structure(file_path, ST, if_start, if_end, condition),
    (ST = "if_statement" ; ST = "unknown"),
    contains(condition, "None"), if_start >= f_start, if_start <= f_end.
.decl RaiseT(file_path: symbol, line: number)
RaiseT(file_path, line) :- source_line(file_path, line, c), contains(c, "raise TypeError").
.decl RaiseV(file_path: symbol, line: number)
RaiseV(file_path, line) :- source_line(file_path, line, c), contains(c, "raise ValueError").
.decl HasBoth(file_path: symbol, class_name: symbol, if_start: number, if_end: number, condition: symbol)
HasBoth(file_path, class_name, if_start, if_end, condition) :-
    IfNone(file_path, class_name, if_start, if_end, condition),
    RaiseT(file_path, l1), l1 >= if_start, l1 <= if_end,
    RaiseV(file_path, l2), l2 >= if_start, l2 <= if_end.
.output HasBoth
\end{minted}
\end{stepbox}
\caption{Complete final query, before semantic validation.}
\label{fig:dataloc-query-repaired}
\end{subfigure}
\caption{Initial and complete final queries from the Astropy interaction, with original identifiers and constraints.}
\label{fig:dataloc-query-repair}
\end{figure}
\end{landscape}

\paragraph*{Localization result.}
After semantic validation, the final query returns the following two candidates:
\begin{center}
\footnotesize
\setstretch{1.1}
\setlength{\tabcolsep}{4pt}
\begin{tabular}{@{}lllll@{}}
\toprule
Repository-relative path & Class & Start & End & Condition \\
\midrule
\texttt{astropy/timeseries/sampled.py} & \texttt{TimeSeries} & 107 & 130 & \texttt{time\_start IsNot None} \\
\texttt{astropy/wcs/wcs.py} & \texttt{WCS} & 391 & 527 & \texttt{header Is None} \\
\bottomrule
\end{tabular}
\end{center}
The LLM retrieves the corresponding source code and checks each candidate against the original request. It excludes \texttt{WCS} because both exception types occur within the same \texttt{else} subtree, rather than in different branches of the conditional. \texttt{TimeSeries} satisfies the branch constraint, yielding the final location:\par\smallskip
\noindent\textbf{Final localization:} \texttt{astropy/timeseries/sampled.py:TimeSeries}
\endgroup
